\setcounter{equation}{0}
\section{Mathematical Basics}

Qubits are conventionally represented as Ket vectors in the \emph{Dirac (Bra-Ket) notation}, which is common in quantum physics texts. The notation defines the Ket column vector, and its row vector transpose the Bra. The Ket vector $\vec{x}$ is denoted by $\ket{\phi}$ = \( \left( \begin{array}{ccc} x_1 \\ x_2 \end{array}\right)\) , and Bra vector $\vec{y}$ is denoted $\bra{\psi}$ = \( \left( \begin{array}{ccc} y_1 & y_2 \end{array}\right)\). In quantum computation, the expression $\ket{0}$ represents quantum zero, and $\ket{1}$ represents quantum one. The state of a qubit can then be represented by the expression $\alpha\ket{0}$+$\beta\ket{1}$, where $\alpha$ and $\beta$ are the amplitudes of zero and one, respectively. The probability of seeing $\ket{0}$ in this state is $\left|\alpha\right|^2$, and $\ket{1}$ is $\left|\beta\right|^2$. As long as our system is coherent, the sum of these probabilities must be 1. So, in this example, $\left|\alpha\right|^2$ + $\left|\beta\right|^2$ = 1. Similarly, an n-qubit register will have $2^n$ amplitudes, each determining the probability that the binary number corresponding to it during measurement, and the sum of their squares will be 1.

There is also an alternative, matrix notation for the quantum computation. The group of amplitudes that describe the state of a quantum register is always a unit vector. According to convention, a qubit with state $\ket{0}$, which is guaranteed to read 0 when measured, is represented by the vector 

\begin{center}
\( \left( \begin{array}{ccc} 1 \\ 0 \end{array}\right)\), 
\end{center}
a qubit with state $\ket{1}$ is represented by the vector 

\begin{center}
\( \left( \begin{array}{ccc} 0 \\ 1 \end{array}\right)\) 
\end{center}
(these are called the basis states). So the amplitude of seeing $\ket{0}$ is given in the first row, and seeing $\ket{1}$ is given in the second row. An arbitrary qubit state is then represented by the vector 

\begin{center}
\( \left( \begin{array}{ccc} \alpha \\ \beta \end{array}\right)\), 
\end{center}
where $\alpha$ and $\beta$ are complex numbers and 

\begin{center}
$\left|\alpha\right|^2$ + $\left|\beta\right|^2$ = 1.
\end{center}
The state of a register consisting of n qubits is represented by the tensor product of the individual states of the qubits in it.

For two matrices \emph{A} and \emph{B}, where \emph{A} is \emph{n} by \emph{m}, \emph{A}$\bigotimes$\emph{B} equals :

\( \left( \begin{array}{ccc} a_{11} & ... & a_{1m} \\ ... & ... & ... \\ a_{n1} & ... & a_{nm} \end{array}\right)\)$\bigotimes$=\( \left( \begin{array}{ccc} a_{11}B & ... & a_{1m}B \\ ... & ... & ... \\ a_{n1}B & ... & a_{nm}B \end{array}\right)\)

So, the tensor product of two vectors $\vec{a}$ and $\vec{b}$ equals :

\begin{center}
\( \left( \begin{array}{ccc} a_1 \\ ... \\ a_n \end{array}\right)\)$\bigotimes$\( \left( \begin{array}{ccc} b_1 \\ ... \\ b_m \end{array}\right)\)=\( \left( \begin{array}{ccc} a_1b_1 \\ ... \\ a_1b_m \\ a_2b_1 \\ ... \\ a_2b_m \\ ... \\ ... \\ a_nb_1 \\ ... \\ a_nb_m \end{array}\right)\)
\end{center}
So a register with two qubits where the bits are \( \left( \begin{array}{ccc} a_1 \\  a_2 \end{array}\right)\) and \( \left( \begin{array}{ccc} b_1 \\  b_2 \end{array}\right)\), respectively can be represented as 


\( \left( \begin{array}{ccc} a_1b_1 \\  a_2b_1 \\ a_1b_2 \\  a_2b_2\end{array}\right)\) where, of course

\begin{center}
$\left|a_1b_1\right|^2$ + $\left|a_1b_2\right|^2$ + $\left|a_2b_1\right|^2$ + $\left|a_2b_2\right|^2$= 1.
\end{center}
A quantum program which transforms a register of $n$ qubits from an input state to an output state can be represented as $2^n$ by $2^n$ matrix, which, when multiplied by the input state vector, gives the output state vector. Because Quantum Computation should be information lossless (ie reversible), this matrix should be a unitary matrix.
The combined effect of several single-qubit programs which run parallelly on the qubits in a register can be represented as the tensor product of their individual matrices.
